\section{Áreas de aplicación práctica de los Agentes}

El apéndice del artículo discute en detalle dos áreas de aplicación más prometedoras para los Agentes, las cuales comparten las siguientes características: requieren la combinación de diálogo y acción, cuentan con criterios de éxito definidos, permiten formar ciclos de retroalimentación e integran una supervisión humana significativa.

\subsection{Soporte al cliente (Customer Support)}

El soporte al cliente es un escenario natural para los Agentes:

\begin{itemize}
    \item La interacción de soporte en sí misma es un flujo de diálogo que, simultáneamente, requiere acceso a información externa y la ejecución de operaciones.
    \item Se pueden integrar herramientas para extraer datos de clientes, historiales de pedidos y artículos de bases de conocimiento.
    \item Operaciones como reembolsos o actualización de tickets de trabajo pueden automatizarse mediante programación.
    \item El éxito puede medirse claramente mediante soluciones validadas por el usuario.
\end{itemize}

\begin{knowledgebox}{Validación del modelo de negocio}
Varias empresas han validado esta dirección mediante un modelo de negocio basado en \textbf{cobranza por solución exitosa} (usage-based pricing) —es decir, solo se cobra cuando el Agente resuelve con éxito un problema del usuario—. Este modelo de negocio es en sí mismo la señal más fuerte de las capacidades del Agente.
\end{knowledgebox}

\subsection{Agentes de programación (Coding Agents)}

El ámbito del desarrollo de software muestra el enorme potencial de los Agentes LLM, cuya capacidad ha evolucionado desde la autocompletación de código hasta la resolución autónoma de problemas. Los Agentes de programación son particularmente efectivos debido a:

\begin{itemize}
    \item Las soluciones de código pueden validarse mediante \textbf{pruebas automatizadas}.
    \item El Agente puede iterar utilizando los resultados de las pruebas como retroalimentación.
    \item El espacio de problemas está bien definido y estructurado.
    \item La calidad de la salida puede medirse objetivamente.
\end{itemize}

\begin{warningbox}{Las pruebas automatizadas no equivalen a una verificación completa}
Aunque las pruebas automatizadas pueden validar la corrección funcional del código, la \textbf{revisión humana} sigue siendo indispensable: se debe asegurar que la solución cumpla con los requisitos más amplios del sistema, las normas de estilo de código, los requisitos de rendimiento y los estándares de seguridad. Aprobar las pruebas es solo el umbral mínimo, no todo lo que constituye garantía de calidad.
\end{warningbox}

\subsection{Resumen de la sección}
El soporte al cliente y la programación son actualmente las áreas de aplicación más maduras para los Agentes. El primero se ajusta a un patrón de interacción conversacional más acciones, mientras que el segundo se beneficia de la testabilidad del código y la estructuración del espacio de problemas. Ambos requieren supervisión humana para garantizar una calidad de nivel superior.
